%
% 1-bernoulli.tex
%
% (c) 2026 Prof Dr Andreas Müller
%
\section{Integration rationaler Funktionen nach Bernoulli
\label{buch:ratint:section:bernoulli}}
\kopfrechts{Integration rationaler Funktionen nach Bernoulli}%

%
% Problemstellung
%
\subsection{Problemstellung}

%
% Stammfunktionen für Partialbrüche
%
\subsection{Stammfunktionen für Partialbrüche}

%
% Rationale Funktionen mit reellen Koeffizienten
%
\subsection{Rationale Funktionen mit reellen Koeffizienten}
Der Fundamentalsatz der Algebra garantiert eine Faktorisierung 
eines Polynoms mit komplexen Koeffizienten in Linearfaktoren
der Form $z-\alpha$, wobei $\alpha$ eine Nullstelle des Polynoms
ist.
Dies ist jedoch nur möglich, wenn wenn komplexe Nullstellen zugelassen
werden.
Es ist jedoch einfach Polynome anzugeben, die über den reellen Zahlen
nicht in Linearfaktoren zerlegt werden können.
Das Polynom $p(x)=x^2+1$ hat keine reellen Nullstellen.
Über den reellen Zahlen ist $p(x)$ ein irreduzibles Polynom.
Über den komplexen Zahlen ist jedoch $p(x)=(x+i)(x-i)$ eine
Faktorisierung in Linearfaktoren.

%
% Faktorisierung in lineare und quadratische Faktoren
%
\subsubsection{Faktorisierung in lineare und quadratische Faktoren}
Für ein reelles Polynom 
$p(z) = a_0 + a_1z + \dots + a_{n-1}z^{n-1} + a_nz^n$
kann man also eine Faktorisierung in Linearfaktoren
nicht erwarten.
Betrachtet man das Polynom jedoch als ein Polynom mit komplexen
Koeffizienten, dann garantiert der Fundamentalsatz der Algebra eine
Faktorisierung
\[
p(z)
=
a_n
(z-\alpha_1) 
(z-\alpha_2) 
\cdot
\ldots
\cdot
(z-\alpha_n),
\]
wobei die $\alpha_k$ Nullstellen des Polynoms sind.
Es gilt also $p(\alpha_k)=0$ für alle $k$.

Das Polynom $p(x)$ hat reelle Koeffizienten, die $\bar{a}_k=a_k$
für alle $k$ erfüllen.
Daher kann man die Bedinung $p(\alpha_k)=0$ komplex konjugieren
und bekommt
\begin{align*}
0
=
\overline{p(\alpha_k)}
&=
\bar{a}_0
+
\bar{a}_1 \bar{\alpha}_k
+
\bar{a}_2 \bar{\alpha}_k^2
+
\ldots
+
\bar{a}_{n-1} \bar{\alpha}_k^{n-1}
+
\bar{a}_{n} \bar{\alpha}_k^{n}
\\
&=
a_0
+
a_1 \bar{\alpha}_k
+
a_2 \bar{\alpha}_k^2
+
\ldots
+
a_{n-1} \bar{\alpha}_k^{n-1}
+
a_{n} \bar{\alpha}_k^{n}
\\
&=
p(\bar{\alpha}_k).
\end{align*}
Es folgt, dass $\bar{\alpha}_k$ ebenfalls eine Nullstelle ist.
Falls $\alpha_k=\bar{\alpha}_k$ ist, ist $\alpha_k$ eine reelle
Nullstelle.
Falls $\alpha_k\ne \bar{\alpha}_k$ ist, dann ist $\bar{\alpha}_k$
eine der anderen Nullstellen.
Die Nullstellen eines reellen Polynoms sind als entweder reell,
oder wenn sie nicht reell sind, dann treten sie in Paaren von
konjugiert komplexen Nullstellen auf.

Ein paar von konjugiert komplexen Nullstellen führt auf den Faktor
\begin{align}
(z-\alpha)(z-\bar{\alpha})
&=
z^2 -(\alpha+\bar{\alpha}) z + \alpha\bar{\alpha}
\label{buch:ratint:bernoulli:reell:faktoren}
\intertext{Die Summe konjugiert komplexer Zahlen ist
$(a+bi)+(a-bi) = 2a $, also das Doppelte des Realteils.
Daher ist der zu berechnende Faktor}
&=
z^2 -2\operatorname{Re}(\alpha)\, z + |\alpha|^2
\notag
\end{align}
Sowohl $\operatorname{Re}\alpha$ und $|\alpha|^2$ sind reell.
Indem man die Paare konjugiert komplexer Zahlen zu quadratischen
Faktoren wie in \eqref{buch:ratint:bernoulli:reell:faktoren}
zusammenfasst, findet man eine Faktorisierung des Polynoms in
der Form
\begin{equation}
p(z)
=
a_n
(z-\alpha_1)
\cdot\ldots\cdot
(z-\alpha_m)
\cdot
(z^2 + \beta_1 z + \gamma_1)
\cdot \ldots \cdot
(z^2 + \beta_l z + \gamma_l),
\label{buch:ratint:bernoulli:eqn:reellefaktorisierung}
\end{equation}
wobei die $\alpha_1,\dots,\alpha_m$ die reellen Nullstellen sind
und die Faktoren $z^2 +\beta_k z +\gamma_k$ Paare von
konjugiert komplexen, aber nicht reellen Lösungen haben.
Damit die rechte Seite den richtigen Grad hat, muss $n=m+2l$ sein.
\label{buch:ratint:bernoulli:eqn:reellefaktorisierung}

%
% Reelle Partialbruchzerlegung
%
\subsubsection{Reelle Partialbruchzerlegung}
In der Zerlegung
\eqref{buch:ratint:bernoulli:eqn:reellefaktorisierung}
können einzelne Faktoren auch mehrfach vorkommen.
Die Konstruktion der Partialbruchzerlegung produziert
daher für jeden Faktor Partialnenner in Potenzen nicht
grösser als in der Faktorisierung.
Dies ist im folgenden Satz zusammengefasst.

\begin{satz}
\label{buch:ratint:bernoulli:satz:reelle-partialbruchzerlegung}
Ist der Nenner $q(z)$ der rationalen Funktion $f(z)=p(z)/q(z)$ 
ein Polynom mit einer Faktorisierung der Form
\begin{equation}
q(z)
=
a_n
(z-\alpha_1)^{s_1}
\cdots
(z-\alpha_m)^{s_m}
(z^2+\beta_1 z + \gamma_1)^{r_1}
\cdots
(z^2+\beta_l z + \gamma_l)^{r_l}
\end{equation}
dann gibt es eine eindeutig bestimmte Partialbruchzerlegung der Form
\begin{equation}
\renewcommand{\arraycolsep}{2pt}
%\renewcommand{\arraystretch}{1.7}
\begin{array}{rcclcccl}
f(z)
&=&&
\displaystyle
\frac{A_1^1    }{ z-\alpha_1       }
&+& \ldots &+&
\displaystyle
\frac{A_1^{r_1}}{(z-\alpha_1)^{s_1}}
\\
&&&\hspace*{11pt}\vdots&&&&\hspace*{17pt}\vdots
\\[2pt]
&&+&
\displaystyle
\frac{A_m^{1}  }{ z-\alpha_m       }
&+& \ldots &+&
\displaystyle
\frac{A_m^{s_m}}{(z-\alpha_m)^{s_m}}
\\[10pt]
&&+&
\displaystyle
\frac{B_1^1    z + C_1^1    }{ z+\beta_1 z + \gamma_1       }
&+& \ldots &+&
\displaystyle
\frac{B_1^{r_1}z + C_1^{r_1}}{(z+\beta_1 z + \gamma_1)^{r_1}}
\\
&&&\hspace*{24pt}\vdots&&&&\hspace*{29pt}\vdots
\\[2pt]
&&+&
\displaystyle
\frac{B_l^1    z + C_1^l    }{ z+\beta_l z + \gamma_l       }
&+& \ldots &+&
\displaystyle
\frac{B_l^{r_l}z + C_l^{r_1}}{(z+\beta_l z + \gamma_l)^{r_l}}
\end{array}
\label{buch:ratint:bernoulli:eqn:reelle-partialbruchzerlegung}
\end{equation}
Darin sind die Zahlen $A_i^k$, $B_i^k$ und $C_i^k$ reelle Zahlen.
\end{satz}

\begin{proof}
Die Form der Partialbruchzerlegung ist aus der früheren Diskussion,
die zu Satz~\ref{buch:analytisch:ratioinal:satz:partialbruchzerlegung}
geführt hat, klar.
Es muss höchstens noch gezeigt werden, dass diese Zerlegung auch
eindeutig ist.

Offenbar sind $s_i$ Zahlen $A_i^k$ zu bestimmen, insgesamt also
$s_1+\dots+s_m$ Unbekannte für Terme zu den Linearfaktoren.
Zu den quadratischen Faktoren sind je $s_i$ Unbekannte
$B_i^k$ und $C_i^k$ zu bestimmen, also insgesamt
$2r_1+\dots+2r_l$ Unbekannte.
Die Gesamtzahl der Unbekannten ist $s_1+\dots+s_m+2r_1+\dots+2r_l=n$.


Multipliziert man die Gleichung
\eqref{buch:ratint:bernoulli:eqn:reelle-partialbruchzerlegung}
mit dem Nenner $q(z)$, dann bleibt auf der linken Seite das Polynom
$p(z)$ und auf der rechten ein Polynom, dessen Koeffizienten
Linearkombinationen der Unbekannten $A_i^k$, $B_i^k$ und $C_i^k$
sind.
Das Zählerpolynom hat Grad $<n$, es enthält also höchstens $n$
Koeffizienten.
Koeffizientenvergleich liefert dann insgesamt $n$ lineare Gleichungen
für die Unbekannten.
Das entstehende lineare Gleichungssystem hat also $n$ Gleichungen für
$n$ Unbekannte.
\end{proof}

Die reelle Partialbruchzerlegung erlaubt die Berechnung einer
Stammfunktion durch Integration der einzelnen Terme.
Für die linearen Terme ist bereits die Stammfunktion
\[
\int\frac{dz}{z-\alpha}
=
\log (z-\alpha)
\]
bekannt.
Die Stammfunktion für die Terme zu den quadratischen Faktoren
des Nenners wird im nächsten Abschnitt bestimmt.

%
% Stammfunktionen für quadratische Faktoren
%
\subsubsection{Stammfunktionen für quadratische Faktoren}
Wir betrachten jetzt einen quadratischen Nenner der Form
\[
z^2 + \beta z + \gamma
=
(z-\alpha)(z-\bar{\alpha}),
\]
wobei $\beta = 2\operatorname{Re}\alpha$ und $\gamma=|\alpha|^2$ ist.

Als erstes berechnen wir das Integral eines Terms, in dessen Nenner
der quadratische Faktor einmal vorkommt, mit Hilfe der komplexen
Partialbruchzerlegung, die die Form
\begin{align*}
\frac{Bz+C}{z^2+\beta z+\gamma}
&=
\frac{A_1}{z-\alpha}
+
\frac{A_2}{z-\bar{\alpha}}.
\intertext{Durch Multiplizieren mit dem Nenner $(z-\alpha)(z-\bar{\alpha})$
entsteht}
Bz+C
&=
A_1(z-\bar{\alpha})
+
A_2(z-\alpha)
=
(A_1+A_2)z-(A_1\bar{\alpha}+A_2\alpha).
\end{align*}
Durch Koeffizientenvergleich entsteht das lineare Gleichungssystem
\[
\renewcommand{\arraycolsep}{2pt}
\begin{array}{rcrcr}
             A_1 &+&        A_2 &=&  B\\
\bar{\alpha} A_1 &+& \alpha A_2 &=& -C\rlap{,}
\end{array}
\]
Das Gleichungssystem ist am einfachsten mit der Carmerschen Regel
zu lösen:
\begin{align*}
A_1
&=
\frac{
\left|\,\begin{matrix*}[r]
     B       &    1   \\
    -C       & \alpha
\end{matrix*}\,\right|
}{
\left|\,\begin{matrix*}[r]
     1       &    1   \\
\bar{\alpha} & \alpha
\end{matrix*}\,\right|
}
=
\frac{B\alpha+C}{2\operatorname{Im}\alpha}
&
A_2
&=
\frac{
\left|\,\begin{matrix*}[r]
     1       &    B   \\
\bar{\alpha} &   -C  
\end{matrix*}\,\right|
}{
\left|\,\begin{matrix*}[r]
     1       &    1   \\
\bar{\alpha} & \alpha
\end{matrix*}\,\right|
}
=
-
\frac{B\bar{\alpha}+C}{2\operatorname{Im}\alpha}.
\end{align*}
Wir wissen bereits, dass der Realteil von $\alpha$ durch
$\operatorname{Re}\alpha=\frac12 \beta$ ausgedrückt werden kann.
Wegen $\gamma=|\alpha|^2$ und $\beta=2\alpha$ kann man
auch den Imaginärteil von $\alpha$ durch $\beta$ und $\gamma$ ausdrücken,
es ist nämlich
\begin{align*}
(\operatorname{Im}\alpha)^2 &= |\alpha|^2 - (\operatorname{Re}\alpha)^2
\\
\operatorname{Im}\alpha &= \frac12\!\sqrt{4\gamma^2 - \beta^2}.
\end{align*}
Damit kann die Lösung auch
\begin{align*}
A_1
&=
\frac{B\alpha + C}{\!\sqrt{4\gamma^2-\beta^2}}
&&&
A_2
&=
-
\frac{B\bar{\alpha} + C}{\!\sqrt{4\gamma^2-\beta^2}}
\\
&=
\frac{B\operatorname{Re}\alpha + C}{\!\sqrt{4\gamma^2-\beta^2}}
+
\frac{iB\operatorname{Im}\alpha}{2\operatorname{Im}\alpha}
&&&
&=
-
\frac{B\operatorname{Re}\alpha + C}{\!\sqrt{4\gamma^2-\beta^2}}
+
\frac{iB\operatorname{Im}\alpha}{2\operatorname{Im}\alpha}
\\
&=
\frac{B\beta + 2C}{2\!\sqrt{4\gamma^2-\beta^2}}
+
\frac{iB}2
&&&
&=
-
\frac{B\beta + 2C}{2\!\sqrt{4\gamma^2-\beta^2}}
+
\frac{iB}2
\end{align*}

Die Partialbruchzerlegung ist also
\[
\frac{Bz+C}{(z-\alpha)(z-\bar{\alpha})}
=
\frac{1}{2\operatorname{Im}\alpha}
\biggl(
\frac{B\alpha +C}{z-\alpha}
-
\frac{B\bar{\alpha} +C}{z-\bar{\alpha}}
\biggr)
\]


In dieser Form kann die Stammfunktion jetzt mit Logarithmusfunktionen
in der Form
\begin{align*}
\int
\frac{Bz+C}{(z-\alpha)(z-\bar{\alpha})}
\,dz
&=
\frac{1}{2\operatorname{Im}\alpha}
\biggl(
(B\alpha+C)
\log(z-\alpha)
-
(B\bar{\alpha}+C)
\log(z-\bar{\alpha})
\biggr)
\\
&=
\frac{1}{2\operatorname{Im}\alpha}
(B\operatorname{Re}\alpha+C)
\log\frac{z-\alpha}{z-\bar{\alpha}}
\\
&\qquad
+
\frac{1}{2\operatorname{Im}\alpha}
\biggl(
iB\operatorname{Im}\alpha
\log(z-\alpha)
+
iB\operatorname{Im}\alpha
\log(z-\bar{\alpha})
\biggr)
\\
&=
\frac{1}{2\operatorname{Im}\alpha}
(B\operatorname{Re}\alpha+C)
\log\frac{z-\alpha}{z-\bar{\alpha}}
+
\frac{iB}{2}
\log\bigl((z-\alpha)(z-\bar{\alpha})\bigr)
)
\\
&=
\frac{B\operatorname{Re}\alpha+C}{2\operatorname{Im}\alpha}
\log
\frac{z-\alpha}{z-\bar{\alpha}}
+
\frac{iB}2
\log(z^2+\beta z+\gamma)
\end{align*}
dargestellt werden.

%
% Stammfunktion für Potenzen von quadratischen Faktoren
%
\subsubsection{Stammfunktion für Potenzen von quadratischen Faktoren}

